\section{前言}

\subsection{研究背景与意义}

本章用于介绍论文研究问题的现实背景、学术价值与应用场景。建议从网络空间安全领域的具体问题出发，说明该问题在真实环境中的影响范围、现有解决手段的局限性，以及继续开展研究的必要性。若选题来自科研项目、企业实践或竞赛课题，可在本节中简要说明来源，但不宜写成项目汇报式语言。

\subsection{国内外研究现状}

本节建议围绕论文主题梳理已有研究进展，可按时间线、技术路线或研究对象进行组织。写作时应使用自己的语言概括相关工作，避免直接复制文献表述；对于关键方法、重要结论或基础定义，应在正文中给出规范引用，例如 \cite{knuth:1984,lesk:1977}。正式写作时，请将模板中的示例引用替换为与课题直接相关的真实文献。

\subsection{本文主要工作与章节安排}

本节可概述论文完成的主要工作，例如问题建模、方案设计、系统实现、实验评估和结果分析等。建议使用一段完整文字说明贡献点，并在最后按章节介绍论文结构，例如“第 2 章介绍预备知识与相关工作，第 3 章介绍方案设计与实现，第 4 章介绍实验与分析，第 5 章总结全文并展望未来工作”。
